\begin{enumerate}
	\item 
	\begin{enumerate}
	\item
	\begin{description}
		\item Avec la formule A, il faudrait payer : $4 \times 45~\text{\euro} = 180~\text{\euro}$.
		\item Avec la formule B, il faudrait payer : $90~\text{\euro} + 4 \times 25,50~\text{\euro} = 192~\text{\euro}$.
		\item Avec la formule C, il faudrait payer $300~\text{\euro}$.
	\end{description}
	
	\medskip
	Si les 4 classes effectuent une visite alors la formule A est la plus avantageuse.
	
	\item Avec 8 visites, le prix à payer serait : 
	\begin{description}
		\item avec la formule A : $8 \times 45~\text{\euro} = 360~\text{\euro}$ ; 
		\item avec la formule B, il faudrait payer : $90~\text{\euro} + 8 \times 25,50~\text{\euro} = 294~\text{\euro}$ ;
		\item avec la formule C, il faudrait payer $300~\text{\euro}$.
	\end{description}
	
	\medskip
	Si les 4 classes effectuent 2 visites, soit 8 visites au total, alors la formule B est la plus avantageuse.
	\end{enumerate}
	
	\item Si on note $x$ le nombre de visites de classes, alors le prix à payer en euros sera de $90+25,5x$ pour la formule B. Il nous faut donc résoudre l'inéquation suivante.
	
	\medskip
	 $90+25,5x > 300$
	 
	 \medskip
	 $25,5x > 210$
	 
	 \medskip
	 $x > \dfrac{210}{25,5}$
	
	\medskip
	$\dfrac{210}{25,5}\approx8,2$ donc, à partir de 9 visites, la formule C est plus économique que la formule B.
	
	\item Pour la formule A, le prix est proportionnel au nombre de visites donc sa représentation graphique est la droite qui passe par l'origine.
	
	\medskip
	Pour la formule B, le prix est donné par l'expression $90+25,5x$. Sa représentation graphique est donc la droite d'ordonnée à l'origine 90.
	
	\medskip
	Pour la formule C, le prix est constant et toujours égal à 300~\euro{}, sa représentation graphique est donc la droite horizontale avec tous les points d'ordonnée 300.
	
	\medskip
	On lit sur le graphique que la droite qui représente la formule B est au-dessous de celle de la formule A à partir de la 5e entrée.\\
	Donc, la formule B est plus économique que la formule A à partir de 5 entrées.
	
	\item On a vu à la question 1. b) que c'était la formule B la plus avantageuse avec un prix de 294~\euro{}. La mairie paie 15~\% du prix total, cela revient donc à une réduction de 15~\% du prix à payer, il faut donc multiplier le prix par : $1 - \dfrac{15}{100} = 0,85$.
	
	$294~\text{\euro{}}\times0,85=249,90~\text{\euro}$
	
	\medskip
	L'école doit prévoir de payer 249,90~\euro.
\end{enumerate}